\subsection{Introduction \quad / \quad \textthai{บทนำ}}
\begin{paracol}{2}
For centuries, scientists have used differential equations to describe the laws of nature. Whether it is the flow of heat, the movement of fluids, or the evolution of quantum wavefunctions, these equations are the bedrock of physics. Traditionally, we solve these using mesh-based numerical methods like Finite Element Analysis (FEA) or Finite Difference Methods (FDM). However, these methods often struggle with complex geometries and high-dimensional spaces.

Physics-Informed Neural Networks\index{PINNs@PINNs (PINNs)} (PINNs) represent a paradigm shift. Instead of discretizing space into a grid, we use a continuous neural network to represent the solution. By embedding the physical laws directly into the loss function, the network learns to satisfy the differential equations inherently.

\switchcolumn

\begin{thai}
เป็นเวลาหลายศตวรรษที่นักวิทยาศาสตร์ใช้สมการเชิงอนุพันธ์เพื่ออธิบายกฎของธรรมชาติ ไม่ว่าจะเป็นการไหลของความร้อน การเคลื่อนที่ของของไหล หรือวิวัฒนาการของคุณสมบัติทางควอนตัม สมการเหล่านี้คือรากฐานของฟิสิกส์ ตามธรรมเนียมเดิมเราแก้สมการเหล่านี้โดยใช้วิธีเชิงตัวเลขแบบอิงตาข่าย (Mesh-based) เช่น Finite Element Analysis (FEA) หรือ Finite Difference Methods (FDM) อย่างไรก็ตาม วิธีการเหล่านี้มักประสบปัญหาเมื่อเจอกับเรขาคณิตที่ซับซ้อนและพื้นที่ที่มีมิติสูง

Physics-Informed Neural Networks (PINNs) เป็นการเปลี่ยนกระบวนทัศน์ แทนที่จะแบ่งพื้นที่ออกเป็นตาราง (Grid) เรากลับใช้โครงข่ายประสาทเทียมแบบต่อเนื่องเพื่อแสดงผลเฉลย การฝังกฎทางฟิสิกส์ลงในฟังก์ชันการสูญเสีย (Loss Function) โดยตรง ทำให้โครงข่ายเรียนรู้ที่จะปฏิบัติตามสมการเชิงอนุพันธ์โดยสัญชาตญาณ
\end{thai}
\end{paracol}
